\scanpage{181}
% archive sheet number: 180
Considérons à nouveau une classe abélienne \deleted{\unsure{de}}
$\unsure{\mathcal K|\mathscr K|\mathcal R}$ dont les trucs soient tous de longueur finie.
Soit $I$ l'ensemble des classes de trucs indécomposables dans $\mathcal K$, soit $Z(I)$ le groupe libre
engendré par $I$. C'est un groupe ordonné $\unsure{\Gamma|F}$, et
d'après Remak-Schmidt, $\Gamma^+$ (l'ens. des
éléments positifs) est canoniquement isomorphe à l'ensemble des classes de trucs
$\in\mathcal K$ modulo isomorphisme, la somme dans
$\Gamma^+$ correspondant à l'opération de somme
directe. Si $\mathcal K$ est une classe de \deleted{\unsure{modules}}, \inserted{\unsure{vectoriels|modules} sur $k$ de dim finie sur $k$ avec structures suppl.}
et si on suppose que $k$ soit muni d'une
structure \unsure{$\mathcal K$|de $K$} privilégiée, alors $I$ a un
élément distingué $e$ ; si dans $K$, on a
une involution $\vee$ (le dual d'un \unsure{vectoriel|variable|module}
muni d'une structure $K$ est muni
d'une structure \unsure{d$K$|$K$}) alors $I$ est muni
d'une involution $i\unsure{-|\to}\check i$, et la
compatibilité de l'opération $\nu$ avec la structure
privilégiée de $k$ s'exprime par $\check e=e$.
\inserted{$\unsure{V|\nu}$ et $e$ \unsure{définissent|déterminent} involution et $e$ dans $\Gamma$}
\unsure{Noter|Notons} que sur $I$, il y a une fonction
privilégiée $n(i)=\dim M_i$, qui se
prolonge par linéarité à $\Gamma=Z(I)$ ;
